% Part: second-order-logic
% Chapter: syntax-and-semantics
% Section: satisfaction

\documentclass[../../../include/open-logic-section]{subfiles}

\begin{document}

\olfileid{sol}{syn}{sat}

\olsection{满足关系}

\begin{explain}
为了定义二阶公式的满足关系 $\Sat{M}{!A}[s]$，我们必须扩展定义以涵盖二阶变元。
二阶逻辑中结构的概念与一阶逻辑中的相同。
唯一的区别在于变元指派~$s$：它现在不仅需要为一阶变元提供值，还需要为二阶变元提供值。
\end{explain}

\begin{defn}[变元指派]
结构~$\Struct{M}$ 的一个\emph{变元指派}~$s$ 是一个函数，它将每个
\begin{enumerate}
\item 对象变元~$\Obj{v_i}$ 映射到~$\Domain M$ 的一个元素，即 $s(\Obj{v_i}) \in \Domain{M}$；
\item $n$元关系变元~$\Obj{V_i^n}$ 映射到~$\Domain{M}$ 上的一个 $n$ 元关系，即 $s(\Obj{V_i^n}) \subseteq \Domain{M}^n$；
\item $n$元函数变元~$\Obj{u_i^n}$ 映射为从 $\Domain{M}$ 到 $\Domain{M}$ 的一个 $n$ 元函数，即 $s(\Obj{u_i^n})\colon \Domain{M}^n \to \Domain{M}$。
\end{enumerate}
\end{defn}

\begin{explain}
结构为每个常项符号和函数符号指定一个值，而二阶变元指派则为每个对象变元和函数变元分别指派对象和函数。
两者结合，使我们可以为每个项指派一个值。
\end{explain}

\begin{defn}[项的值]
如果 $t$ 是语言~$\Lang L$ 的一个项，$\Struct M$ 是~$\Lang L$ 的一个结构，且 $s$ 是~$\Struct M$ 的一个变元指派，
则 \emph{值}~$\Value{t}{M}[s]$ 的定义与一阶项相同，但需增加以下子句：
\begin{quote}
\indcase{t}{\Atom{u}{t_1, \ldots, t_n}}{
\[
\Value{\indfrm}{M}[s] = s(u)(\Value{t_1}{M}[s], \ldots,
\Value{t_n}{M}[s]).
\]}
\end{quote}
\end{defn}

\begin{defn}[$x$-变体]
如果 $s$ 是结构~$\Struct M$ 的一个变元指派，那么任何与 $s$ 至多在对 $x$ 的赋值上有所不同的、$\Struct M$ 的变元指派 $s'$，称为 $s$ 的一个\emph{$x$-变体}。
若 $s'$ 是 $s$ 的一个 $x$-变体，我们记作 $\varAssign{s'}{s}{x}$。（对二阶变元 $X$ 或 $u$ 类似。）
\end{defn}

\begin{defn}
  如果 $s$ 是结构~$\Struct M$ 的一个变元指派且 $m \in \Domain{M}$，
  那么指派~$\Subst{s}{m}{x}$ 是如下定义的变元指派：
  \[\Subst{s}{m}{y} = \begin{cases}
    m & \text{若 } y \ident x\\
    s(y) & \text{否则},
  \end{cases}\]
  如果 $X$ 是一个 $n$ 元关系变元且 $M \subseteq \Domain{M}^n$，那么 $\Subst{s}{M}{X}$ 是如下定义的变元指派：
  \[\Subst{s}{M}{y} = \begin{cases}
    M & \text{若 } y \ident X\\
    s(y) & \text{否则}. 
  \end{cases}\]
  如果 $u$ 是一个 $n$ 元函数变元且 $f\colon \Domain{M}^n
  \to \Domain{M}$，那么 $\Subst{s}{f}{u}$ 是如下定义的变元指派：
  \[\Subst{s}{f}{y} = \begin{cases}
    f & \text{若 } y \ident u\\
    s(y) & \text{否则}.
  \end{cases}\]
  在上述各情形中，$y$ 可以是任意一阶或二阶变元。
\end{defn}

\begin{defn}[满足]
对于二阶公式~$!A$，满足的定义与 \olref[fol][syn][sat]{defn:satisfaction} 类似，但需补充：
\begin{enumerate}
\item \indcase{!A}{\Atom{X^n}{t_1, \dots, t_n}}{$\Sat{M}{\indfrm}[s]$ 当且仅当 $\langle \Value{t_1}{M}[s], \dots, \Value{t_n}{M}[s] \rangle \in
  s(X^n)$。}
\tagitem{prvAll}{%
  \indcase{!A}{\lforall[X][!B]}{$\Sat{M}{\indfrm}[s]$ 当且仅当对每个 $M \subseteq \Domain{M}^n$，$\Sat{M}{!B}[\Subst{s}{M}{X}]$。}}{}

\tagitem{prvEx}{%
  \indcase{!A}{\lexists[X][!B]}{$\Sat{M}{\indfrm}[s]$ 当且仅当存在至少一个 $M \subseteq \Domain{M}^n$，使得 $\Sat{M}{!B}[\Subst{s}{M}{X}]$。}}{}

\tagitem{prvAll}{%
  \indcase{!A}{\lforall[u][!B]}{$\Sat{M}{\indfrm}[s]$ 当且仅当对每个 $f\colon \Domain{M}^n \to \Domain{M}$，
  $\Sat{M}{!B}[\Subst{s}{f}{u}]$。}}{}

\tagitem{prvEx}{%
  \indcase{!A}{\lexists[u][!B]}{$\Sat{M}{\indfrm}[s]$ 当且仅当存在至少一个 $f\colon \Domain{M}^n \to \Domain{M}$，使得 $\Sat{M}{!B}[\Subst{s}{f}{u}]$。}}{}
\end{enumerate}
\end{defn}

\begin{ex}
  考虑公式 $\lforall[z][(\Atom{X}{z} \liff \lnot
    \Atom{Y}{z})]$。它不包含二阶量词，但包含二阶变元 $X$ 和~$Y$（这里视为一元的）。对应的一阶语句
  $\lforall[z][(\Atom{P}{z} \liff \lnot \Atom{R}{z})]$ 说的是：任何归属于~$P$ 的解释的对象都不归属于~$R$ 的解释，反之亦然。在结构中，谓词符号~$P$ 的解释由解释函数~$\Assign{P}{M}$ 给出。但对于像 $X$ 和~$Y$ 这样的二阶变元，解释不是由结构本身，而是由一个变元指派提供的。由于该二阶公式不是语句（它包含自由变元 $X$ 和~$Y$），它只能相对于结构~$\Struct{M}$ 连同变元指派~$s$ 被满足。

  $\Sat{M}{\lforall[z][(\Atom{X}{z} \liff \lnot \Atom{Y}{z})]}[s]$ 当且仅当~$s(X)$ 的元素都不是~$s(Y)$ 的元素，反之亦然，即 $s(Y) = \Domain{M} \setminus s(X)$。例如，考虑 $\Domain{M} = \{1, 2, 3\}$。由于不涉及谓词符号、函数符号或常项符号，~$\Struct{M}$ 的论域是唯一有关的因素。若令 $s_1(X) = \{1, 2\}$
  和 $s_1(Y) = \{3\}$，我们有 $\Sat{M}{\lforall[z][(\Atom{X}{z}
      \liff \lnot \Atom{Y}{z})]}[s_1]$。

  相反，若 $s_2(X) = \{1, 2\}$ 且 $s_2(Y) = \{2, 3\}$，则 $\Sat/{M}{\lforall[z][(\Atom{X}{z} \liff \lnot
      \Atom{Y}{z})]}[s_2]$。这是因为 $\Sat{M}{\Atom{X}{z}}[\Subst{s_2}{2}{z}]$（因为 $2 \in \Subst{s_2}{2}{z}(X)$），但 $\Sat/{M}{\lnot \Atom{Y}{z}}[\Subst{s_2}{2}{z}]$（因为 $2 \in
  \Subst{s_2}{2}{z}(Y)$）。
\end{ex}

\begin{ex}
  $\Sat{M}{\lexists[Y][(\lexists[y][\Atom{Y}{y}] \land
  \lforall[z][(\Atom{X}{z} \liff \lnot \Atom{Y}{z})])]}[s]$ 当且仅当存在 $N \subseteq \Domain{M}$ 使得 $\Sat{M}{(\lexists[y][\Atom{Y}{y}] \land \lforall[z][(\Atom{X}{z}
  \liff \lnot \Atom{Y}{z})])}[\Subst{s}{N}{Y}]$。且对任意 $N \neq \emptyset$（从而 $\Sat{M}{\lexists[y][\Atom{Y}{y}]}[\Subst{s}{N}{Y}]$）以及如前例那样 $N = \Domain{M} \setminus s(X)$，该情况均成立。换句话说，$\Sat{M}{\lexists[Y][(\lexists[y][\Atom{Y}{y}] \land
  \lforall[z][(\Atom{X}{z} \liff \lnot \Atom{Y}{z})])]}[s]$ 当且仅当 $\Domain{M} \setminus s(X)$ 非空，即 $s(X) \neq \Domain{M}$。因此，例如当 $\Domain{M}
  = \{1, 2, 3\}$ 且 $s(X) = \{1, 2\}$ 时，该公式被满足；但当 $s(X) = \{1, 2, 3\}
  = \Domain{M}$ 时则不被满足。

  由于只要 $s(X) = \Domain{M}$ 该公式就不被满足，语句
  \[
  \lforall[X][\lexists[Y][(\lexists[y][\Atom{Y}{y}] \land
      \lforall[z][(\Atom{X}{z} \liff \lnot \Atom{Y}{z})])]]
  \]
  就永远不会被满足：对任何结构~$\Struct{M}$，指派~$s(X) = \Domain{M}$ 都将使该语句为假。另一方面，语句
  \[
  \lexists[X][\lexists[Y][(\lexists[y][\Atom{Y}{y}] \land
      \lforall[z][(\Atom{X}{z} \liff \lnot \Atom{Y}{z})])]]
  \]
  相对于任何指派~$s$ 都被满足，因为我们总可以找到 $M \subseteq \Domain{M}$ 但 $M \neq \Domain{M}$（例如 $M = \emptyset$）。
\end{ex}

\begin{ex}
  二阶语句~$\lforall[X][\lforall[y][X(y)]]$ 说的是每个一元关系，即每个性质，对每个对象都成立。这显然永远不会为真，因为在每个~$\Struct{M}$ 中，对于满足 $s(X) = \emptyset$ 且 $s(y) = a \in
  \Domain{M}$ 的变元指派~$s$，我们有 $\Sat/{M}{X(y)}[s]$。这意味着 $!A \lif
  \lforall[X][\lforall[y][X(y)]]$ 在二阶逻辑中等价于 $\lnot !A$，即：$\Sat{M}{!A \lif
  \lforall[X][\lforall[y][X(y)]]}$ 当且仅当 $\Sat{M}{\lnot !A}$。换言之，在二阶逻辑中，我们可以用 $\lforall$ 和~$\lif$ 来定义 $\lnot$。
\end{ex}

\begin{prob}
  证明在二阶逻辑中，$\lforall$ 和~$\lif$ 可以定义其他联结词：
  \begin{enumerate}
    \item 证明在二阶逻辑中，$!A \land !B$ 等价于 $\lforall[X][(!A \lif (!B \lif \lforall[x][X(x)])\lif
    \lforall[x][X(x)])]$。
    \item 找出一个仅使用 $\lforall$ 和 $\lif$ 的二阶公式，使其等价于 $!A \lor !B$。
  \end{enumerate}
\end{prob}

\end{document}